\section*{第 \ref{ch:b2:planet-formation-meteorites} 章——行星形成与陨石}

\begin{solution}{exo:b2:planet-formation-meteorites:1}
气体：\qty{2200}{g.cm^{-2}}（位于 \qty{1}{au} 处）；在 \qty{5.2}{au} 处，
$2200\times5.2^{-3/2} = 2200/11.86 = \qty{186}{g.cm^{-2}}$。固体：
$33\times0.33 = \qty{10.9}{g.cm^{-2}}$（位于 \qty{1}{au} 处）；
$33/11.86 = \qty{2.78}{g.cm^{-2}}$（位于 \qty{5.2}{au} 处）。在 \qty{5.2}{au} 处，气体
$186/2.78 = 67$ 倍于固体的质量。
\end{solution}

\begin{solution}{exo:b2:planet-formation-meteorites:2}
铝（\qty{1653}{K}，超难熔）、镁（\qty{1336}{K}，常见），
铁（\qty{1334}{K}，常见）、钠（\qty{958}{K}，中等挥发性）、硫
（\qty{664}{K}，挥发性）、以水冰形式存在的氧（\qty{180}{K}，高挥发性）。
CAI（\omterm{def:b2:planet-formation-meteorites:chondrule}{富钙铝包体}）是最早的凝结物，富集铝（以及钙、钛）；
相对于镁和铝，地球中的挥发性元素（钠、硫和水）
明显贫乏。
\end{solution}

\begin{solution}{exo:b2:planet-formation-meteorites:3}
目击陨落：普通\omterm{def:b2:planet-formation-meteorites:chondrule}{球粒陨石} $856/1107 = \qty{77.3}{\percent}$，\omterm{def:b2:planet-formation-meteorites:classes}{铁陨石}
$49/1107 = \qty{4.4}{\percent}$。\omterm{def:b2:planet-formation-meteorites:fall-find}{发现型陨石}：$38647/44534 = \qty{86.8}{\percent}$ 与
$1020/44534 = \qty{2.3}{\percent}$。南极\omterm{def:b2:planet-formation-meteorites:fall-find}{发现型陨石}：$22096/44534 =
\qty{49.6}{\percent}$。目击陨落：每一次被观测到的坠落，无论属于哪一类都会被收集，
而\omterm{def:b2:planet-formation-meteorites:fall-find}{发现型陨石}则取决于\omterm{def:b2:planet-formation-meteorites:fall-find}{陨石}在地表的保存程度，
外观如何，以及人们在何处搜寻；其中一半出自同一类
地点。
\end{solution}

\begin{solution}{exo:b2:planet-formation-meteorites:4}
$\Delta t = 1.035\ln(\num{5.23e-5}/\num{1.0e-5}) = 1.035\ln 5.23 =
1.035\times1.654 = \qty{1.71}{Ma}$，即在 CAI 之后。
\end{solution}

\begin{solution}{exo:b2:planet-formation-meteorites:5}
$4\pi\Sigma_1 r_1^2 = 4\pi\times\num{2.2e4}\times(\num{1.496e11})^2 =
\qty{6.19e27}{kg}$。在 \qty{30}{au} 以内：$\times\sqrt{30} = 5.48$，则
\qty{3.39e28}{kg} $= 0.017$ 个太阳质量；在 \qty{100}{au} 以内：$\times10$，
\qty{6.19e28}{kg} $= 0.031$ 个太阳质量。两者均位于观测范围之内，
且更接近其上限：该星云是普通的盘。
\end{solution}

\begin{solution}{exo:b2:planet-formation-meteorites:6}
$r = (120/T)^{7/3}$ au：对于水，$(120/182)^{7/3} = 0.659^{7/3} =
\qty{0.38}{au}$；对于甲烷，$(120/41)^{7/3} = 2.93^{7/3} = \qty{12.3}{au}$。
直接受阳光照射的尘粒位于盘面表层；盘的中平面，
受上方尘埃遮蔽而温度更低，因此冰能在更靠内的位置留存。
\omterm{def:b2:planet-formation-meteorites:planetesimal}{星子}形成于中平面，因此中平面的估计适用，前提是
盘为被动盘；吸积加热会使\omterm{def:b2:planet-formation-meteorites:snow-line}{雪线}外移。
\end{solution}

\begin{solution}{exo:b2:planet-formation-meteorites:7}
在 \qty{1}{au} 处：$r^2 = \num{2.238e22}$，$\Sigma_p = \qty{109}{kg.m^{-2}}$，于是
$2\pi b r^2\Sigma_p = 31.42\times\num{2.238e22}\times109 = \num{7.66e25}$;
$(\num{7.66e25})^{3/2} = \num{6.71e38}$，再除以
$(3\times\num{1.989e30})^{1/2} = \num{2.443e15}$：\qty{2.75e23}{kg}
$= 0.046$ 个地球质量。在 \qty{5.2}{au} 处：$r^2 = \num{6.05e23}$，
$\Sigma_p = \qty{27.8}{kg.m^{-2}}$, $2\pi b r^2\Sigma_p = \num{5.29e26}$,
$M = \num{1.22e40}/\num{2.443e15} = \qty{4.98e24}{kg} = 0.83$ 个地球质量。
若固体量增至 \num{3.5} 倍，$M$ 按 $3.5^{3/2} = 6.55$ 增长，约为 5.5
个地球质量。
\end{solution}

\begin{solution}{exo:b2:planet-formation-meteorites:8}
每千克含 $\num{1.01e19}$ 个原子，因此
$\num{1.01e19}\times\num{5.00e-13} = \qty{5.05e6}{J.kg^{-1}}$。温度
升高将为 $\num{5.05e6}/837 = \qty{6030}{K}$。由 \num{250} 升温至
\qty{1700}{K} 需要 $837\times1450 = \qty{1.21e6}{J.kg^{-1}}$：铝-26
提供的能量约为该值的四倍。
\end{solution}

\begin{solution}{exo:b2:planet-formation-meteorites:9}
$\num{5.07e10}/\num{7.49e7} = 677$。铁-60 的平均寿命为
$2.62/\ln2 = \qty{3.78}{Ma}$。每克的初始功率：
$\num{5.07e10}/1.035 = \num{4.90e10}$ 与 $\num{7.49e7}/3.78 = \num{1.98e7}$ 尔格
每克每百万年，比值约为 \num{2470}：铝-26 占据绝对优势。
\end{solution}

\begin{solution}{exo:b2:planet-formation-meteorites:10}
功率相等的条件为 $\num{4.90e10}\,e^{-t/1.035} = \num{1.98e7}\,e^{-t/3.78}$，
即 $t(1/1.035 - 1/3.78) = \ln(\num{4.90e10}/\num{1.98e7}) = \ln 2475 = 7.81$；
其中 $0.966 - 0.265 = 0.702$ 按每百万年计，$t = \qty{11.1}{Ma}$。此时
铁-60 剩余 $\num{7.49e7}\,e^{-11.1/3.78} = \num{7.49e7}\times0.053 =
\qty{4.0e6}{erg.g^{-1}}$，即 \qty{400}{J.kg^{-1}}，少了数千倍，
不足以熔化任何物质。
\end{solution}

\begin{solution}{exo:b2:planet-formation-meteorites:11}
$\tau = 1.035\times\num{3.156e13} = \qty{3.267e13}{s}$，且
$H_0\tau = \qty{6.21e6}{J.kg^{-1}}$。对于 \qty{1200}{K}：
$t_a^* = 1.035\ln\bigl(\num{6.21e6}/(837\times950)\bigr) = 1.035\ln7.81 =
\qty{2.13}{Ma}$；对于 \qty{1700}{K}：$1.035\ln\bigl(\num{6.21e6}/(837\times1450)\bigr)
= 1.035\ln5.12 = \qty{1.69}{Ma}$。将 $H_0$ 减半，两者各减去 $\tau\ln2 =
\qty{0.72}{Ma}$，因为 $t_a^* = \tau\ln H_0 + \tau\ln(\tau/c(T_c -
T_0))$，且只有第一项发生变化：偏移量恰为一个半衰期，
与临界温度无关。
\end{solution}

\begin{solution}{exo:b2:planet-formation-meteorites:12}
$\kappa = 2.1/(3300\times837) = \qty{7.60e-7}{m^2.s^{-1}}$，且
$\sqrt{\kappa\tau} = \sqrt{\num{7.60e-7}\times\num{3.267e13}} = \qty{4980}{m}$，
约 \qty{5}{km}。直径为 \qty{100}{km} 的天体，其中心深度为 \qty{50}{km}，
是该距离的十倍：在铝-26 衰变期间其热量无法散失，因此该
命题适用。直径为 \qty{2}{km} 的天体，其散热速率
与其产热速率相当。熔化需要的半径为 \qty{5}{km} 的数倍，
即直径约 \qtyrange{10}{20}{km} 或更大，正如模型所示。
\end{solution}

\begin{solution}{pb:b2:planet-formation-meteorites:1}
\textbf{1.} 目击陨落：其坠落被人看到，次日清晨被拾获。
它未经风化，并计入抵达地球的物质的无偏样本。

\textbf{2.} $856/1107 = 0.77$：约四分之三的机会。

\textbf{3.} 球粒直径 \qty{0.3}{mm}，含 \qty{9}{\percent} 金属及橄榄石，
其中铁橄榄石占 \qty{18}{\percent}：属 H 型\omterm{def:b2:planet-formation-meteorites:chondrule}{球粒陨石}（金属含量 \qty{10}{\percent}，Fa 16--20）。

\textbf{4.} 球粒边界模糊、矿物成分均一，说明经历了热变质：
类型为 5 至 6 型。普通\omterm{def:b2:planet-formation-meteorites:chondrule}{球粒陨石}所达到的温度至多约 \qty{1220}{K}，即
\qty{947}{\degreeCelsius}。

\textbf{5.} 普通\omterm{def:b2:planet-formation-meteorites:chondrule}{球粒陨石}中 CAI 仅占 \qtyrange{0.1}{1}{\percent}（按体积计）。
该岩石由球粒和基质组成，它们在 CAI 之后聚集而成，
CAI 是最古老的固体，任何含有它们的岩石都更年轻。

\textbf{6.} 每种矿物在形成时都获得相同的镁，以及与其铝含量成比例的铝-26；
此后该铝-26 已衰变为镁-26，
因此矿物中每个镁原子所对应的铝越多，
其获得的镁-26 也越多。

\textbf{7.} $(0.139590 - 0.139500)/(15.0 - 0.05) = \num{0.000090}/14.95 =
\num{6.02e-6}$：球粒形成时的比值 
$\ce{^{26}Al}/\ce{^{27}Al}$。

\textbf{8.} $0.139500 + \num{6.02e-6}\times3.00 = 0.139518$：位于该直线上。

\textbf{9.} $\Delta t = 1.035\ln(\num{5.23e-5}/\num{6.02e-6}) = 1.035\ln8.69 =
1.035\times2.162 = \qty{2.24}{Ma}$。

\textbf{10.} $4567.5 - 2.2 = \qty{4565.3}{\Ma}$。

\textbf{11.} $\num{5.23e-5}\,e^{-10/1.035} = \num{5.23e-5}\times\num{6.4e-5} =
\num{3.3e-9}$，镁-26 含量过少，无法测量。铪-182（半衰期
\qty{8.9}{Ma}）或锰-53（\qty{3.7}{Ma}）适于研究
CAI 之后约一千万年发生的事件。

\textbf{12.} $\kappa = 2.1/(3300\times837) = \qty{7.60e-7}{m^2.s^{-1}}$；
$R^2/\kappa = \num{4e8}/\num{7.60e-7} = \qty{5.26e14}{s} = \qty{16.7}{Ma}$，是铝-26 平均寿命的十六倍：
天体中心保留了其热量。

\textbf{13.} $\num{1.9e-7}\,e^{-2.4/1.035} = \num{1.9e-7}\times0.0984 =
\qty{1.87e-8}{W.kg^{-1}}$。

\textbf{14.} $\num{1.87e-8}\times\num{3.267e13} = \qty{6.11e5}{J.kg^{-1}}$。

\textbf{15.} $\Delta T = \num{6.11e5}/837 = \qty{730}{K}$；峰值
温度为 $250 + 730 = \qty{980}{K}$。

\textbf{16.} 高于 \qty{750}{K}，因此岩石发生烧结并经历变质；低于
\qty{1200}{K}，因此金属和硫化物均未熔化。该\omterm{def:b2:planet-formation-meteorites:fall-find}{陨石}为经受加热的\omterm{def:b2:planet-formation-meteorites:chondrule}{球粒陨石}，
\omterm{def:b2:planet-formation-meteorites:type}{岩石学类型}较高，低于第 4 题中的 \qty{1220}{K}：二者一致。

\textbf{17.} $1.035\ln\bigl(\num{6.21e6}/(837\times950)\bigr) = \qty{2.13}{Ma}$：
约早 \qty{0.3}{Ma}。

\textbf{18.} 温度的上升量为 $\Delta T(1 - e^{-(t - t_a)/\tau})$；达到 \qty{90}{\percent}
时需满足 $e^{-(t - t_a)/\tau} = 0.1$，$t - t_a = 1.035\ln10 = \qty{2.38}{Ma}$，即
CAI 之后约 \qty{4.8}{Ma}。

\textbf{19.} $\Delta T = 7420\,e^{-0.8/1.035} = 7420\times0.462 = \qty{3430}{K}$，
峰值为 \qty{3680}{K}，远高于硅酸盐熔融时的 \qty{1700}{K}。
天体发生了熔化（熔化热使其温度维持在较低水平），
并分异为金属地核，以及硅酸盐地幔与地壳。

\textbf{20.} $e^{(2.4 - 0.8)/1.035} = e^{1.546} = 4.7$ 倍。

\textbf{21.} 锌和铅在 \qty{727}{K} 附近凝结，属中等挥发性元素；镁
（\qty{1336}{K}）和铝（\qty{1653}{K}）凝结得早得多。该\omterm{def:b2:planet-formation-meteorites:fall-find}{陨石}所在区域的固体
在气体冷却到足以使全部锌和铅凝结之前，
固体便已聚集。

\textbf{22.} $33\times0.33\times2.5^{-3/2} = 10.89/3.95 = \qty{2.76}{g.cm^{-2}}$，即
\qty{27.6}{kg.m^{-2}}。

\textbf{23.} $r^2 = (3.74\times10^{11})^2 = \num{1.40e23}$；
$2\pi\times5\times\num{1.40e23}\times27.6 = \num{1.21e26}$；
$(\num{1.21e26})^{3/2}/\num{2.443e15} = \qty{5.45e23}{kg}$，即 $0.091$ 个地球质量。
母天体的质量为 $\tfrac43\pi(\num{2e4})^3\times3300 = \qty{1.1e17}{kg}$，比前者小
五百万倍：它是\omterm{def:b2:planet-formation-meteorites:planetesimal}{星子}，而非\omterm{def:b2:planet-formation-meteorites:planetesimal}{行星胚胎}。

\textbf{24.} 母天体中心的峰值温度约为
\textbf{\qty[text-series-to-math=true]{707}{\degreeCelsius}}（$980 - 273$）：足以使\omterm{def:b2:planet-formation-meteorites:chondrule}{球粒陨石}重结晶，
但不足以使其金属熔化。
\end{solution}
